\section{Related work}\label{sec:related}

\paragraph{Misconceptions and logic errors of novice programmers.}
Reviews of introductory programming describe misconceptions about variables, control flow, parameter passing and the notional machine that executes a program \citep{sorva2013notional,qian2017misconceptions}. Error studies distinguish slips from errors that follow from a wrong model of the language or task \citep{altadmri2015compilations,hermans2021brain}. \citet{ettles2018logic} classified logic errors of CS1 students into algorithmic errors, misinterpretations of the task and misconceptions, and found misconceptions to be the most frequent and the hardest to resolve. Misconception-driven feedback can help students when the misconception is known in advance \citep{gusukuma2018misconception}, and interactive instruments that ask students to predict program behaviour can elicit misconceptions directly \citep{lu2024smol}. These studies motivate automatic grouping of errors, but they also show that a program alone rarely reveals which belief produced it; we therefore evaluate program-level mechanisms and keep cognitive claims out of scope.

\paragraph{Clustering student programs.}
OverCode clusters correct solutions by normalised code \citep{glassman2015overcode}, and Clara clusters correct solutions by dynamic equivalence to repair incorrect ones \citep{gulwani2018clara}. For incorrect submissions, \citet{kaleeswaran2016semisupervised} cluster by repair strategy so that an instructor fixes one representative per cluster; MistakeBrowser clusters incorrect submissions by the code transformations that students applied when they fixed their own programs \citep{head2017mistakebrowser}; TraceDiff contrasts the traces of a buggy program and its closest fix \citep{suzuki2017tracediff}. Program embeddings learned from code and execution have been used to propagate feedback \citep{piech2015embeddings}, and \citet{shi2021misconception} clustered embeddings of incorrect block-based programs to help experts discover misconceptions. For C, InvAASTCluster combines dynamic invariants with abstracted syntax trees \citep{orvalho2025invaast}, and MENTOR couples clustering, formula-based fault localisation and LLM repair \citep{orvalho2026mentor}. These systems use repairs or correct solutions as part of the method. We use the student's own later repair only to construct evaluation labels, and ask how well representations that are available \emph{before} any repair recover them.

\paragraph{Mining misconceptions with learned models.}
Subtree-attention models localise logical errors from correctness labels alone \citep{hoq2025sann}, pattern-based knowledge components are extracted from student code with representation learning and clustering \citep{edm2026patternkc}, and LLMs discover misconceptions from sets of student code samples in McMining, whose benchmark injects known misconceptions into Python code \citep{alhossami2026mcmining}. Recent classroom tools cluster LLM-identified errors for instructors \citep{sigcse2026realtime}, ground tutoring in known misconceptions \citep{sigcse2026tutor}, or refine automatically detected signals with expert review \citep{aruleba2026ukicer}. PyMETA shows that diagnosis accuracy depends strongly on how the gold label is defined: prompted models that label the first execution error well agree far less with expert repair paths \citep{li2026pymeta}. Our study shares this concern and makes the label definition explicit, verifiable and reproducible.

\paragraph{Benchmarks, mutants and repairs.}
C-Pack-IPAs provides C90 submissions of 25 assignments with anonymised student identifiers, tests and historical grading logs \citep{orvalho2024cpack}; ITSP is an earlier benchmark of buggy and corrected C programs \citep{yi2017itsp}. Mutation analysis injects controlled faults; mutant detection correlates with real-fault detection in software testing \citep{just2014mutants}, and conceptual mutants built from students' wrong examples expose specification misconceptions \citep{prasad2024conceptual}. Delta debugging isolates a minimal failure-inducing difference between a failing and a passing configuration \citep{zeller2002ddmin}. We combine these ideas: we replay the corpus, reduce each student's own repair to a verified minimal change, and pair these labels with controlled single-mechanism mutants.

\paragraph{Rule induction.}
The Inductive Learning Algorithm (ILA) extracts production rules class by class from attribute--value examples \citep{tolun1998ila}; ILA-2 adds a penalty factor that tolerates noisy data \citep{tolun1999ila2}. ILA is sometimes listed alongside clustering methods, but it is supervised rule induction, and we use it as such: to state IF--THEN rules that map execution evidence to mechanisms and to measure how well such rules hold on students not used to induce them.
